%/*% pdflatex -file-line-error -interaction nonstopmode %

\documentclass{article}
\usepackage{graphicx}
\usepackage{amsmath}
\usepackage{amsfonts}
\usepackage{bookmark}
\usepackage{booktabs}
\usepackage{exercises}
\graphicspath{ {./images/test} }
\usepackage{indentfirst}
\usepackage{amsthm}
\usepackage{yhmath}
\usepackage{placeins}
\newcommand*{\oldneg}{\mathord{\sim}}
\newcommand\showdiv[1]{\overline{\smash{)}#1}}
\newcommand\Mydiv[1]{%
\smash{\raise.3ex\hbox{$\big)$}}\mkern-8mu
        \overline{\enspace\strut#1}}
\setlength{\tabcolsep}{5pt}
\theoremstyle{definition}
\newtheorem{definition}{Definition}

\setlength{\parindent}{.5in}
\setlength{\parskip}{12pt}
\renewcommand{\arraystretch}{2}

\title{3.1 Answers}
\begin{document}
\maketitle

\section*{1}
In the expression $\sqrt[4]{10,000}$, the number $4$ is called the index of the radical.

\section*{2}
(a): Suppose that we are taking the real fourth roots of $a$. If $a$ is positive, this means that there are two solutions. If $a$ is negative, this means that there 
are not real solutions. This is because any real number raised to an even power (4) is a positive number, and therefore cannot equal a negative number ($a$). If 
$a$ is zero, than there is one root, namely $0$. \newline
(b): Suppose that we are taking the real fifth roots of some number $a$. If $a$ is positive, this means there is exactly one real root, namely a positive number. 
If $a$ is negative, then there is exactly one real root, namely a negative number. This is because, if a number is raised to an odd power (5), it keeps whatever sign 
it originally had. If $a=0$, then there is one root, namely $0$.

\section*{5}
D

\section*{6}
B

\section*{9}
\[ \sqrt{x^5}=(x^5)^\frac{1}{2}=x^\frac{5}{2} \]

\section*{10}
\[ \sqrt[5]{(y+2)^2}=((y+2)^2)^\frac{1}{5}=(y+2)^\frac{2}{5} \]

\section*{13}
\[ (2t)^\frac{1}{3}=\sqrt[3]{2t} \]

\section*{14}
\[ x^\frac{7}{6}=(x^7)^\frac{1}{6}=\sqrt[6]{x^7} \]

\section*{19}
Because $n=4$ is even and $a=-16$ is negative, we know that there are no real $n$th roots in this situation.

\section*{20}
Because $n=5$ is odd and $a=-32$ is negative, we know that there is one negative real $n$th root. It is $-2$.

\section*{31}
\[ 64^{-\frac{2}{3}}=(\sqrt[3]{64})^{-2}=4^{-2}=\frac{1}{16} \]

\section*{33}
\[ 128^\frac{5}{7}=\sqrt[7]{128}^5=2^5=32 \]
C

\section*{45}
\[ \frac{1}{(-17)^\frac{3}{5}}=\frac{1}{\sqrt[5]{-17}^3}=-\frac{1}{5.47}\approx -\frac{1}{5.47}\]

\section*{46}
\[ 27^\frac{3}{5}\approx 7.22 \]
\[ 5^\frac{3}{2}\approx 11.18 \]
\[ \sqrt[3]{81}\approx 4.33 \]
\[ (\sqrt[3]{2})^8\approx 6.35 \]
B

\section*{47}
Because $a$ is positive and $n$ is even, there are two $4$th roots. We know that $3^4=81$ and $4^4=256$. Similarly $(-3)^4=81$ and $(-4)^4=256$. Thus, we know that the 
positive root lies between $3$ and $4$, and the negative real root lies between $-4$ and $-3$.

\section*{49}
The error here is that, because $n=4$ is an even number and $a=81$ is a positive number, there are 2 $4$th roots, one positive, one negative. The answer only gives 
the positive root. To fix the answer, we can add $x=-3$ to be the second $4$th root. 

\section*{57}
\[ (x+10)^5=70 \]
\[ x+10=\sqrt[5]{70}\approx 2.34 \]
\[ x=-7.66 \]

\section*{58}
\[ x^6-34=181 \]
\[ x^6=215 \]
\[ x=\pm\sqrt[6]{215}\approx\pm 2.45 \]

\section*{59}
(a): From the graph, we can see that the range is $y\ge 0$, and for every $y\ge 0$, there are two $x$'s such that $y=x^n$, $\pm\sqrt[n]{a}$, 
which corresponds exactly to the Key Concept 
box in the case where $n$ is an even integer.\newline
(b): This graph has a shape such that the range is $-\infty <y<+\infty$, and for every $y$, there is one $x$ such that $y=x^n$, a positive $x$ is $y>0$, and a negative 
$x$ if $y<0$, which matches the Key Concept box in the case where $n$ is an odd integer.

\section*{66}
Because the mass of the particles that a river can transport is proportional to the sixth power of the speed of the river, \[ m=ks^6 \] where $k$ is some constant. 
If a river normally flows at a speed of 1 meter per second, then the mass of the particles that a river can transport is $m=k1^6=k$. To transport particles 10 times 
as massive, we would have $10k=10m=ks^6$ and thus $s=\sqrt[6]{10}\approx 1.47$ meters per second. 
To transport particles 100 times as massive, we would have $100k=100m=ks^6$ and thus $s=\sqrt[6]{100}\approx 2.15$ meters per second.

\end{document}

